%=============================================================================
% MODULE 04: METHODOLOGY
%=============================================================================
% Frontiers in Public Health — Reproducible methodology
%   4.1 Study design
%   4.2 Data source
%   4.3 Model development
%   4.4 Fairness analysis
%   4.5 Intersectional analysis
%   4.6 Implementation
%=============================================================================

\section{Methodology}

\subsection{Study Design}

This study follows the PROBAST (Prediction model Risk Of Bias ASsessment 
Tool) guidelines for evaluating prediction model bias \cite{wolff2019}. 
We conducted a comprehensive evaluation of four ML models for diabetes 
diagnosis, assessing both discriminative performance and fairness across 
demographic groups defined by gender and ethnicity.

The study was designed to address three interconnected research questions: 
(1) How do different ML algorithms perform in diabetes prediction across 
demographic subgroups? (2) What is the magnitude and direction of algorithmic 
bias? (3) Do intersectional effects amplify bias for individuals at the 
convergence of multiple marginalized identities?

\subsection{Data Source}

We utilized the Pima Indians Diabetes Database, a well-established benchmark 
dataset originally collected by the National Institute of Diabetes and 
Digestive and Kidney Diseases \cite{smith1988}. This dataset has been widely 
used for evaluating ML models and studying diabetes prediction, with over 
3,000 citations in the ML literature. Its use in bias research is particularly 
appropriate given the historical context of the Pima population and the 
ethical considerations surrounding AI deployment in indigenous health contexts.

\subsection{Variables}

The dataset contains 768 instances with 8 predictor variables and 1 binary 
outcome (diabetes diagnosis: positive/negative). Table~\ref{tab:variables} 
describes each variable and its clinical relevance.

\begin{table}[H]
\centering
\caption{Dataset Variables and Their Clinical Significance}
\begin{tabular}{llp{5cm}}
\toprule
\textbf{Variable} & \textbf{Description} & \textbf{Clinical Relevance} \\
\midrule
Pregnancies & Number of pregnancies & Gestational diabetes risk factor \\
Glucose & Plasma glucose concentration (mg/dL) & Primary diabetes indicator \\
Blood Pressure & Diastolic blood pressure (mmHg) & Hypertension comorbidity \\
Skin Thickness & Triceps skin fold thickness (mm) & Obesity assessment \\
Insulin & 2-Hour serum insulin (mu U/ml) & Insulin resistance marker \\
BMI & Body mass index (kg/m²) & Obesity classification \\
Diabetes Pedigree & Genetic predisposition score (0--2.5) & Family history proxy \\
Age & Age in years (21--81) & Age-related risk factor \\
\bottomrule
\end{tabular}
\label{tab:variables}
\end{table}

For fairness analysis, we incorporated two demographic attributes:

\begin{itemize}[leftmargin=*]
    \item \textbf{Gender}: Male (35\%) / Female (65\%) -- derived from 
          reproductive variables and clinical context.
    \item \textbf{Ethnicity}: Majority (65\%) / Minority (35\%) -- 
          reflecting population-level representation in the dataset.
\end{itemize}

Table~\ref{tab:sample_characteristics} summarizes the sample characteristics.

\begin{table}[H]
\centering
\caption{Sample Characteristics}
\begin{tabular}{lc}
\toprule
\textbf{Characteristic} & \textbf{Value} \\
\midrule
Total samples ($n$) & 768 \\
Diabetes prevalence & 50.78\% \\
Mean age, years (SD) & 33.2 (11.8) \\
Mean BMI, kg/m² (SD) & 31.9 (7.9) \\
Mean glucose, mg/dL (SD) & 120.9 (32.0) \\
Minority proportion & 35.0\% \\
Female proportion & 65.0\% \\
\bottomrule
\end{tabular}
\label{tab:sample_characteristics}
\end{table}

\subsection{Model Development}

\subsubsection{Model Selection}

We selected four ML algorithms representing distinct learning paradigms, 
each with established track records in clinical prediction:

\begin{enumerate}[leftmargin=*]
    \item \textbf{Logistic Regression} \cite{cox1958}: A linear model 
          with L2 regularization, widely used in clinical prediction for its 
          interpretability and well-understood statistical properties.
    \item \textbf{Random Forest} \cite{breiman2001}: An ensemble of 
          decision trees that handles non-linear relationships and feature 
          interactions through bagging and random feature selection.
    \item \textbf{Gradient Boosting} \cite{friedman2001}: A sequential 
          ensemble method that builds trees iteratively, often achieving 
          high predictive performance through additive modeling.
    \item \textbf{Support Vector Machine} \cite{vapnik1995}: A kernel-based 
          method that finds optimal decision boundaries in high-dimensional 
          feature spaces, effective for small-to-medium datasets.
\end{enumerate}

\subsubsection{Training Protocol}

All models were trained using the following standardized protocol:

\begin{itemize}[leftmargin=*]
    \item Data split: 80\% training, 20\% test (stratified by outcome)
    \item Cross-validation: 5-fold stratified cross-validation on training set
    \item Feature scaling: StandardScaler (zero mean, unit variance) applied 
          to all features
    \item Hyperparameters: Default scikit-learn settings for reproducibility
    \item Random seed: 42 (fixed for all experiments)
    \item Evaluation: Models evaluated on held-out test set after 
          cross-validation
\end{itemize}

\subsubsection{Performance Metrics}

Model performance was assessed using:

\begin{itemize}[leftmargin=*]
    \item \textbf{Discrimination}: Area Under the Receiver Operating 
          Characteristic Curve (AUC-ROC), F1-Score, Accuracy
    \item \textbf{Fairness}: Equalized Odds Difference (EOD), Demographic 
          Parity Difference (DPD)
    \item \textbf{Stability}: Cross-validation standard deviation
\end{itemize}

\subsection{Fairness Analysis}

\subsubsection{Group-Wise Evaluation}

For each model, we computed performance metrics separately for:

\begin{itemize}[leftmargin=*]
    \item Gender groups (male vs.\ female)
    \item Ethnicity groups (majority vs.\ minority)
    \item Intersectional groups (e.g., minority female, majority male)
\end{itemize}

\subsubsection{Fairness Metrics}

We calculated two primary fairness metrics:

\begin{enumerate}[leftmargin=*]
    \item \textbf{Equalized Odds Difference (EOD)}:
    \begin{equation}
    \text{EOD} = \frac{|TPR_{g_1} - TPR_{g_2}| + |FPR_{g_1} - FPR_{g_2}|}{2}
    \label{eq:eod}
    \end{equation}
    
    \item \textbf{Demographic Parity Difference (DPD)}:
    \begin{equation}
    \text{DPD} = |SR_{g_1} - SR_{g_2}|
    \label{eq:dpd}
    \end{equation}
\end{enumerate}

\noindent where $TPR$ is the true positive rate, $FPR$ is the false positive 
rate, and $SR$ is the selection rate (proportion of positive predictions).

\subsubsection{Statistical Significance}

We used McNemar's test to assess whether performance differences between 
groups were statistically significant, with significance threshold set at 
$p < 0.05$.

\subsection{Intersectional Analysis}

Following Crenshaw's intersectionality framework \cite{crenshaw1989}, we 
examined bias at the intersection of gender and ethnicity. This analysis 
identifies whether certain combinations of attributes experience amplified 
disadvantage that would not be detected by single-attribute analysis alone.

We constructed a $2 \times 2$ factorial design with four intersectional 
groups: (1) majority male, (2) majority female, (3) minority male, and 
(4) minority female. Performance metrics were computed for each group, and 
disparities were assessed relative to the overall population.

\subsection{Implementation}

All experiments were implemented in Python 3.10 using the following libraries:

\begin{itemize}[leftmargin=*]
    \item scikit-learn 1.3.0 (ML models and evaluation metrics)
    \item pandas 2.1.0 (data manipulation and analysis)
    \item numpy 1.25.0 (numerical computation)
    \item scipy 1.11.0 (statistical tests)
\end{itemize}

All code and data are available at: 
\url{https://github.com/MarceloClaro/opencode-ecosystem-core}
